\chapter{Magnétostatique}

\newpage

\section{Solénoïde épais à densité de courant variable}

Un bobinage cylindrique d'axe $Oz$, de rayon $a$ et de longueur $L\gg a$, est parcouru par des courants orthoradiaux. La densité de spires croît avec la distance à l'axe, de sorte que la densité volumique de courant s'écrit
\begin{align*}
	\vec{\jmath}(r)=j_0\frac{r}{a}\,\vec{e}_\theta
\end{align*}
On se place loin des extrémités du bobinage, et on prendra $a=10$~cm et $j_0=5,0$~A$\cdot$mm$^{-2}$.

\begin{figure}[!h]
\centering
\includegraphics[width=0.80\linewidth]{magnetostatique1_1.pdf}
\end{figure}

\begin{enumerate}

	\item À l'aide des invariances et des symétries de la distribution, déterminer la direction du champ magnétique et les variables dont il dépend.

	\item En appliquant le théorème d'Ampère au contour $\Gamma$ de la figure, exprimer $B(r)$ en fonction de $B(0)$ pour $r<a$. Que devient ce résultat pour $r>a$ ?

	\item Que peut-on dire du champ à l'extérieur du bobinage ? En déduire $B(0)$, puis $\vec{B}(r)$ en tout point. Tracer l'allure de $B(r)$.

	\item Retrouver ce résultat en décomposant le bobinage en solénoïdes élémentaires emboîtés. Commenter le profil obtenu.

\end{enumerate}

\newpage

\begin{correction}

\begin{enumerate}

	\item \textbf{Invariances.} La distribution est invariante par rotation autour de $Oz$ et, loin des extrémités, par translation selon $z$ : le champ ne dépend donc que de $r$.

	\textbf{Symétries.} Les courants étant orthoradiaux, ils sont contenus dans tout plan $z=\mathrm{cte}$, qui est donc un plan de symétrie de la distribution. Le champ magnétique étant un \textit{pseudo-vecteur}, il est perpendiculaire à un plan de symétrie en tout point de ce plan :
	\begin{align*}
		\boxed{\vec{B}=B(r)\,\vec{e}_z}
	\end{align*}
	\textit{Vérification par l'autre voie : tout plan contenant $Oz$ change $\vec{e}_\theta$ en $-\vec{e}_\theta$, c'est donc un plan d'antisymétrie, dans lequel un pseudo-vecteur est contenu. $\vec{e}_z$ appartient bien à ces plans.}

	\item Le champ étant porté par $\vec{e}_z$, les deux côtés radiaux de $\Gamma$ ne contribuent pas à la circulation. En parcourant $\Gamma$ dans le sens indiqué sur la figure,
	\begin{align*}
		\oint_\Gamma \vec{B}\cdot\dif\vec{l}=h\,B(0)-h\,B(r)
	\end{align*}
	Ce sens correspond, par la règle de la main droite, à une normale orientée selon $+\vec{e}_\theta$ : le courant enlacé vaut
	\begin{align*}
		I_{\mathrm{enl}}=\int_0^h\!\!\int_0^r j_0\frac{r'}{a}\,\dif r'\,\dif z=\frac{h\,j_0r^2}{2a}
	\end{align*}
	et le théorème d'Ampère donne
	\begin{align*}
		\boxed{B(r)=B(0)-\frac{\mu_0j_0r^2}{2a}}\qquad(r<a)
	\end{align*}
	Pour $r>a$, le contour n'enlace plus que la totalité du courant, $I_{\mathrm{enl}}=hj_0a/2$, d'où
	\begin{align*}
		B(r)=B(0)-\frac{\mu_0j_0a}{2}\qquad(r>a)
	\end{align*}
	qui ne dépend plus de $r$ : \textbf{le champ est uniforme à l'extérieur du bobinage}.

	\item Le bobinage est long, mais fini : à grande distance, le champ tend vers zéro. Or il vient d'être établi qu'il est uniforme pour $r>a$, et une constante qui tend vers zéro est nulle. Donc $B(r>a)=0$, ce qui fixe enfin la constante :
	\begin{align*}
		\boxed{B(0)=\frac{\mu_0j_0a}{2}}
		\qquad\text{et}\qquad
		\boxed{\vec{B}(r)=\frac{\mu_0j_0}{2a}\left(a^2-r^2\right)\vec{e}_z\;\;(r<a),
		\qquad \vec{B}=\vec{0}\;\;(r>a)}
	\end{align*}
	\textit{Application numérique : $B(0)=\mu_0j_0a/2=0,31$~T, l'ordre de grandeur d'un électroaimant de laboratoire.}

	\begin{center}
	\includegraphics[width=0.90\linewidth]{magnetostatique1_2.pdf}
	\end{center}

	\textit{Remarque pour le colleur : c'est l'argument du solénoïde infini, et il vaut la peine de le faire expliciter. Pour un cylindre \textit{rigoureusement} infini, le problème n'a pas de solution unique : on peut toujours superposer un champ uniforme $B_0\vec{e}_z$, qui vérifie les équations de Maxwell dans le vide. C'est la finitude du bobinage, donc la décroissance du champ à l'infini, qui lève l'indétermination. Un élève qui répond « $\vec{B}$ est nul dehors, c'est un solénoïde » a le bon réflexe mais pas encore l'argument.}

	\item Une couche comprise entre $r'$ et $r'+\dif r'$ porte le courant surfacique $\dif K=j(r')\dif r'$ : c'est un solénoïde élémentaire, qui crée $\mu_0\dif K\,\vec{e}_z$ à l'intérieur et rien à l'extérieur. Au rayon $r$, seules les couches telles que $r'>r$ contribuent :
	\begin{align*}
		B(r)=\int_r^a\mu_0j_0\frac{r'}{a}\,\dif r'=\frac{\mu_0j_0}{2a}\left(a^2-r^2\right)
	\end{align*}
	ce qui redonne le résultat précédent sans aucune constante d'intégration à fixer.

	Le champ est donc \textbf{maximal sur l'axe et nul au bord}, exactement l'inverse du fil massif parcouru par un courant axial, où $B$ est nul sur l'axe et maximal en surface. La décomposition en solénoïdes le rend évident : plus on s'éloigne de l'axe, moins il reste de couches extérieures pour contribuer.

	\textit{Remarque pour le colleur : sur la réalisation du dispositif. Un courant orthoradial permanent est impossible dans un conducteur massif, puisqu'il faudrait une circulation non nulle de $\vec{E}$ le long d'une ligne de courant fermée, ce qu'interdit le régime stationnaire. La distribution proposée décrit donc un \textbf{bobinage} dont la densité de spires croît linéairement avec $r$, et non un bloc conducteur. Le courant total par unité de longueur vaut $\int_0^a j\,\dif r=j_0a/2=2,5$~kA$\cdot$cm$^{-1}$ : c'est l'équivalent d'un solénoïde mince parcouru par ce courant linéique.}

\end{enumerate}

\end{correction}
